\begin{textsample}{POS Dim 1 – 2012 – Score 20.00 – 2012/t002337.txt}  \label{ex:f1_pos_036}
America ' s public energy conversation boils down to this question : Would you rather die of A ) \textbf{oil} wars, or B ) climate change, or C ) nuclear holocaust, or D ) all of \textbf{the} above?
Oh, I missed one : or E ) none of \textbf{the} above?
That ' s \textbf{the} one we ' re not normally offered.
What if we could make energy do our work without working our undoing?
Could we have fuel without fear?
Could we reinvent fire?
You see, fire made us human ; \textbf{fossil} fuels made us modern.
But now we need a new fire that makes us safe, secure, healthy and durable.
Let ' s see how.
Four-fifths of \textbf{the} world ' s energy still comes from burning each year four cubic miles of \textbf{the} rotted remains of primeval swamp goo.
Those \textbf{fossil} fuels have built our civilization.
They ' ve created our wealth.
They ' ve enriched \textbf{the} lives of billions.
But they also have rising costs to our security, economy, health and environment that are starting to erode, if not outweigh their benefits.
So we need a new fire.
And switching from \textbf{the} old fire to \textbf{the} new fire means changing two big stories about \textbf{oil} and electricity, each of which puts two- fifths of \textbf{the} \textbf{fossil} \textbf{carbon} in \textbf{the} air.
But they ' re really quite distinct.
Less than one percent of our electricity is made from \textbf{oil} — although almost half is made from \textbf{coal}.
Their uses are quite concentrated.
Three-fourths of our \textbf{oil} fuel is \textbf{transportation}.
Three-fourths of our electricity powers buildings.
And \textbf{the} rest of both runs factories.
So very efficient vehicles, buildings and factories save \textbf{oil} and \textbf{coal}, and also natural gas that can displace both of them.
But today ' s energy system is not just inefficient, it is also disconnected, aging, dirty and insecure.
So it needs refurbishment.
By 2050 though, it could become efficient, connected and distributed with elegantly frugal autos, factories and buildings all relying on a modern, secure and resilient electricity system.
We can eliminate our addiction to \textbf{oil} and \textbf{coal} by 2050 and use \textbf{one-third} less natural gas while switching to efficient use and \textbf{renewable} supply.
This could cost, by 2050, five trillion dollars less in net present value, that is expressed as a lump \textbf{sum} today, than business as usual — assuming that \textbf{carbon} \textbf{emissions} and all other hidden or external costs are worth zero — a conservatively low estimate.
Yet this cheaper energy system could support 158 percent bigger U.
S. economy all without needing \textbf{oil} or \textbf{coal}, or for that matter nuclear energy.
Moreover, this transition needs no new inventions and no acts of Congress and no new federal taxes, mandate subsidies or laws and running Washington gridlock.
Let me say that again.
I ' m going to tell you how to get \textbf{the} United States completely off \textbf{oil} and \textbf{coal}, five trillion dollars cheaper with no act of Congress led by business for profit.
In other words, we ' re going to use our most effective institutions — private enterprise co-evolving with civil society and sped by military innovation to go around our least effective institutions.
And whether you care most about profits and jobs and competitive advantage or national security, or environmental stewardship and climate protection and public health, reinventing fire makes sense and makes money.
General Eisenhower reputedly said that enlarging \textbf{the} boundaries of a tough problem makes it soluble by encompassing more options and more synergies.
So in reinventing fire, we integrated all four sectors that use energy — \textbf{transportation}, buildings, industry and electricity — and we integrated four kinds of innovation, not just technology and policy, but also design and business strategy.
Those combinations yield very much more than \textbf{the} \textbf{sum} of \textbf{the} parts, especially in creating deeply disruptive business opportunities.
Oil costs our economy two billion dollars a day, plus another four billion dollars a day in hidden economic and military costs, raising its total cost to over a sixth of GDP.
Our mobility fuel goes three-fifths to automobiles.
So let ' s start by making autos \textbf{oil} free. \textbf{Two-thirds} of \textbf{the} energy it takes to move a typical car is caused by its weight.
And every unit of energy you save at \textbf{the} wheels, by taking out weight or drag, saves seven units in \textbf{the} tank, because you don ' t have to waste six units getting \textbf{the} energy to \textbf{the} wheels.
Unfortunately, over \textbf{the} past quarter century, epidemic obesity has made our two-ton steel cars gain weight twice as fast as we have.
But today, ultralight, ultrastrong materials, like \textbf{carbon} fiber composites, can make dramatic weight savings snowball and can make cars simpler and cheaper to build.
Lighter and more slippery autos need less force to move them, so their engines get smaller.
Indeed, that sort of vehicle fitness then makes electric propulsion affordable because \textbf{the} batteries or fuel cells also get smaller and lighter and cheaper.
So sticker prices will ultimately fall to about \textbf{the} same as today, while \textbf{the} driving cost, even from \textbf{the} start, is very much lower.
So these innovations together can transform automakers from wringing tiny savings out of Victorian engine and seal-stamping technologies to \textbf{the} steeply falling costs of three linked innovations that strongly reenforce each other — namely ultralight materials, making them into structures and electric propulsion. \textbf{The} sales can grow and \textbf{the} prices fall even faster with temporary feebates, that is rebates for efficient new autos paid for by fees on inefficient ones.
And just in \textbf{the} first two years \textbf{the} biggest of Europe ' s five feebate programs has tripled \textbf{the} speed of improving automotive efficiency. \textbf{The} resulting shift to electric autos is going to be as game-changing as shifting from typewriters to \textbf{the} gains in \textbf{computers}.
Of course, \textbf{computers} and electronics are now America ' s biggest industry, while typewriter makers have vanished.
So vehicle fitness opens a new automotive competitive strategy that can double \textbf{the} \textbf{oil} savings over \textbf{the} next 40 years, but then also make electrification affordable, and that displaces \textbf{the} rest of \textbf{the} \textbf{oil}.
America could lead this next automotive revolution.
Currently \textbf{the} leader is Germany.
Last year, Volkswagen announced that by next year they ' ll be producing this \textbf{carbon} fiber plugin hybrid getting 230 miles a gallon.
Also last year, BMW announced this \textbf{carbon} fiber electric car, they said that its \textbf{carbon} fiber is paid for by needing fewer batteries.
And they said,"We do not intend to be a typewriter maker."Audi claimed it ' s going to beat them both by a year.
Seven years ago, an even faster and cheaper American manufacturing technology was used to make this little \textbf{carbon} fiber test part, which doubles as a \textbf{carbon} cap.
In one minute — and you can tell from \textbf{the} sound how immensely stiff and strong it is.
Don ' t worry about dropping it, it ' s tougher than titanium.
Tom Friedman actually whacked it as hard as he could with a sledgehammer without even scuffing it.
But such manufacturing \textbf{techniques} can scale to automotive speed and cost with aerospace performance.
They can save four-fifths of \textbf{the} capital needed to make autos.
They can save lives because this stuff can absorb up to 12 times as much crash energy per \textbf{pound} as steel.
If we made all of our autos this way, it would save \textbf{oil} equivalent to finding one and a half Saudi Arabias, or half an OPEC, by drilling in \textbf{the} Detroit formation, a very prospective play.
And all those mega-barrels under Detroit cost an average of 18 \textbf{bucks} a barrel.
They are all- American, carbon-free and inexhaustible. \textbf{The} same physics and \textbf{the} same business logic also apply to big vehicles.
In \textbf{the} five years ending with 2010, Walmart saved 60 percent of \textbf{the} fuel per ton-mile in its giant fleet of heavy trucks through better logistics and design.
But just \textbf{the} technological savings in heavy trucks can get to \textbf{two-thirds}.
And combined with triple to quintuple efficiency \textbf{airplanes}, now on \textbf{the} drawing board, can save close to a trillion dollars.
Also today ' s military revolution in energy efficiency is going to speed up all of these civilian advances in much \textbf{the} same way that military R\&D has given us \textbf{the} Internet, \textbf{the} Global Positioning System and \textbf{the} jet engine and microchip industries.
As we design and build vehicles better, we can also use them smarter by harnessing four powerful \textbf{techniques} for eliminating needless driving.
Instead of just seeing \textbf{the} travel grow, we can use innovative pricing, charging for road infrastructure by \textbf{the} mile, not by \textbf{the} gallon.
We can use some smart IT to enhance transit and enable car sharing and ride sharing.
We can allow smart and lucrative growth models that help people already be near where they want to be, so they don ' t need to go somewhere else.
And we can use smart IT to make traffic free-flowing.
Together, those things can give us \textbf{the} same or better access with 46 to 84 percent less driving, saving another 0. 4 trillion dollars, plus 0. 3 trillion dollars from using trucks more productively.
So 40 years hence, when you add it all up, a far more mobile U.
S. economy can use no \textbf{oil}.
Saving or displacing barrels for 25 \textbf{bucks} rather than buying them for over a hundred, adds up to a \$4 trillion net saving counting all \textbf{the} hidden costs at zero.
So to get mobility without \textbf{oil}, to phase out \textbf{the} \textbf{oil}, we can get efficient and then switch fuels.
Those 125 to 240 mile-per-gallon-equivalent autos can use any mixture of hydrogen fuel cells, electricity and advanced biofuels. \textbf{The} trucks and planes can realistically use hydrogen or advanced biofuels. \textbf{The} trucks could even use natural gas.
But no vehicles will need \textbf{oil}.
And \textbf{the} most biofuel we might need, just three million barrels a day, can be made \textbf{two-thirds} from waste without displacing any cropland and without harming soil or climate.
Our team speeds up these kinds of \textbf{oil} savings by what we call"institutional acupuncture."We figure out where \textbf{the} business logic is congested and not flowing properly, we stick little needles in it to get it flowing, working with partners like Ford and Walmart and \textbf{the} Pentagon.
And \textbf{the} long transition is already well under way.
In fact, three years ago mainstream analysts were starting to see peak \textbf{oil}, not in supply, but in demand.
And Deutsche Bank even said world \textbf{oil} use could peak around 2016.
In other words, \textbf{oil} is getting uncompetitive even at low prices before it becomes unavailable even at high prices.
But \textbf{the} electrified vehicles don ' t need to burden \textbf{the} electricity grid.
Rather, when smart autos exchange electricity and information through smart buildings with smart grids, they ' re adding to \textbf{the} grid valuable flexibility and storage that help \textbf{the} grid integrate varying solar and wind power.
So \textbf{the} electrified autos make \textbf{the} auto and electricity problems easier to solve together than separately.
And they also converge \textbf{the} \textbf{oil} story with our second big story, saving electricity and then making it differently.
And those twin revolutions in electricity will bring to that sector more numerous and profound and diverse disruptions than any other sector, because we ' ve got 21st century technology and speed colliding head-on with 20th and 19th century institutions, rules and cultures.
Changing how we make electricity gets easier if we need less of it.
Most of it now is wasted and \textbf{the} technologies for saving it keep improving faster than we ' re installing them.
So \textbf{the} unbought efficiency resource keeps getting ever bigger and cheaper.
But as efficiency in buildings and industry starts to grow faster than \textbf{the} economy, America ' s electricity use could actually shrink, even with \textbf{the} little extra use required for those efficient electrified autos.
And we can do this just by reasonably accelerating existing trends.
Over \textbf{the} next 40 years, buildings, which use three-quarters of \textbf{the} electricity, can triple or quadruple their energy productivity, saving 1. 4 trillion dollars, net present value, with a 33 percent internal rate of return or in English, \textbf{the} savings are worth four times what they cost.
And industry can accelerate too, doubling its energy productivity with a 21 percent internal rate of return. \textbf{The} key is a disruptive innovation that we call integrative design that often makes very big energy savings cost less than small or no savings.
That is, it can give you expanding returns, not diminishing returns.
That is how our 2010 retrofit is saving over two-fifths of \textbf{the} energy in \textbf{the} Empire State Building — remanufacturing those six and a half thousand windows on site into super windows that pass light, but reflect heat. plus better lights and office equipment and such cut \textbf{the} maximum cooling load by a third.
And then renovating smaller chillers instead of adding bigger ones saved 17 million dollars of capital cost, which helped pay for \textbf{the} other improvements and reduce \textbf{the} payback to just three years.
Integrative design can also increase energy savings in industry.
Dow ' s billion-dollar efficiency investment has already returned nine billion dollars.
But industry as a whole has another half-trillion dollars of energy still to save.
For example, three-fifths of \textbf{the} world ' s electricity runs motors.
Half of that runs pumps and fans.
And those can all be made more efficient, and \textbf{the} motors that turn them can have their system efficiency roughly doubled by integrating 35 improvements, paying back in about a year.
But first we ought to be capturing bigger, cheaper savings that are normally ignored and are not in \textbf{the} \textbf{textbooks}.
For example, pumps, \textbf{the} biggest use of motors, move liquid through pipes.
But a standard industrial pumping loop was redesigned to use at least 86 percent less energy, not by getting better pumps, but just by replacing long, thin, crooked pipes with \textbf{fat}, short, straight pipes.
This is not about new technology, it ' s just rearranging our metal furniture.
Of course, it also shrinks \textbf{the} pumping equipment and its capital costs.
So what do such savings mean for \textbf{the} electricity that is three-fifths used in motors?
Well, from \textbf{the} \textbf{coal} burned at \textbf{the} power plant through all these compounding losses, only a tenth of \textbf{the} fuel energy actually ends up coming out \textbf{the} pipe as flow.
But now let ' s turn those compounding losses around \textbf{backwards}, and every unit of flow or friction that we save in \textbf{the} pipe saves 10 units of fuel cost, pollution and what Hunter Lovins calls"global weirding"back at \textbf{the} power plant.
And of course, as you go back upstream, \textbf{the} components get smaller and therefore cheaper.
Our team has lately found such snowballing energy savings in more than 30 billion dollars worth of industrial redesigns — everything from data centers and chip fabs to mines and refineries.
Typically our retrofit designs save about 30 to 60 percent of \textbf{the} energy and pay back in a few years, while \textbf{the} new facility designs save 40 to 90-odd percent with generally lower capital cost.
Now needing less electricity would ease and speed \textbf{the} shift to new sources of electricity, chiefly renewables.
China leads their explosive growth and their plummeting cost.
In fact, these solar power \textbf{module} costs have just fallen off \textbf{the} bottom of \textbf{the} chart.
And Germany now has more solar workers than America has steel workers.
Already in about 20 states private installers will come put those cheap solar cells on your roof with no money down and beat your utility bill.
Such unregulated products could ultimately add up to a virtual utility that bypasses your electric company just as your cellphone bypassed your wireline phone company.
And this sort of thing gives utility executives \textbf{the} heebee-jeebees and it gives venture capitalists sweet dreams.
Renewables are no longer a fringe activity.
For each of \textbf{the} past four years half of \textbf{the} world ' s new generating capacity has been \textbf{renewable}, mainly lately in developing countries.
In 2010, renewables other than big hydro, particularly wind and solar cells, got 151 billion dollars of private investment, and they actually surpassed \textbf{the} total installed capacity of nuclear power in \textbf{the} world by adding 60 billion watts in that one year.
That happens to be \textbf{the} same amount of solar cell capacity that \textbf{the} world can now make every year — a number that goes up 60 or 70 percent a year.
In contrast, \textbf{the} net additions of nuclear capacity and \textbf{coal} capacity and \textbf{the} orders behind those keep fading because they cost too much and they have too much financial risk.
In fact in this country, no new nuclear power plant has been able to raise any private construction capital, despite seven years of 100-plus percent subsidies.
So how else could we replace \textbf{the} coal-fired power plants?
Well efficiency and gas can displace them all at just below their operating cost and, combined with renewables, can displace them more than 23 times at less than their replacement cost.
But we only need to replace them once.
We ' re often told though that only \textbf{coal} and nuclear plants can keep \textbf{the} lights on, because they ' re 24 / 7, whereas wind and solar power are variable, and hence supposedly unreliable.
Actually no generator is 24 / 7.
They all break.
And when a big plant goes down, you lose a thousand megawatts in milliseconds, often for weeks or months, often without warning.
That is exactly why we ' ve designed \textbf{the} grid to back up failed plants with working plants.
And in exactly \textbf{the} same way, \textbf{the} grid can handle wind and solar power ' s forecastable variations.
Hourly simulations show that largely or wholly \textbf{renewable} grids can deliver highly reliable power when they ' re forecasted, integrated and diversified by both type and location.
And that ' s true both for continental areas like \textbf{the} U.
S. or Europe and for smaller areas embedded within a larger grid.
That is how, for example, four German states in 2010 were 43 to 52 percent wind powered.
Portugal was 45 percent \textbf{renewable} powered, Denmark 36.
And it ' s how all of Europe can shift to \textbf{renewable} electricity.
In America, our aging, dirty and insecure power system has to be replaced anyway by 2050.
And whatever we replace it with is going to cost about \textbf{the} same, about six trillion dollars at present value — whether we buy more of what we ' ve got or new nuclear and so- called clean \textbf{coal}, or renewables that are more or less centralized.
But those four futures at \textbf{the} same cost differ profoundly in their risks, around national security, fuel, water, finance, technology, climate and health.
For example, our over-centralized grid is very vulnerable to cascading and potentially economy-shattering blackouts caused by bad space weather or other natural disasters or a \textbf{terrorist} attack.
But that blackout risk \textbf{disappears}, and all of \textbf{the} other risks are best managed, with distributed renewables organized into local micro-grids that normally interconnect, but can stand alone at need.
That is, they can disconnect fractally and then reconnect seamlessly.
That approach is exactly what \textbf{the} Pentagon is adopting for its own power supply.
They think they need that ; how about \textbf{the} rest of us that they ' re defending?
We want our stuff to work too.
At about \textbf{the} same cost as business as usual, this would maximize national security, customer choice, entrepreneurial opportunity and innovation.
Together, efficient use and diverse dispersed \textbf{renewable} supply are starting to transform \textbf{the} whole electricity sector.
Traditionally utilities build a lot of giant \textbf{coal} and nuclear plants and a bunch of big gas plants and maybe a little bit of efficiency renewables.
And those utilities were rewarded, as they still are in 34 states, for selling you more electricity.
However, especially where regulators are now instead rewarding cutting your bills, \textbf{the} investments are shifting radically toward efficiency, demand response, cogeneration, renewables and ways to knit them all together reliably with less transmission and little or no bulk electricity storage.
So our energy future is not fate, but choice, and that choice is very flexible.
In 1976, for example, government and industry insisted that \textbf{the} amount of energy needed to make a dollar of GDP could never go down.
And I heretically suggested it could go down several-fold.
Well that ' s what ' s actually happened so far.
It ' s fallen by half.
But with today ' s much better technologies, more mature delivery channels and integrative design, we can do far more and even cheaper.
So to solve \textbf{the} energy problem, we just needed to enlarge it.
And \textbf{the} results may at first seem incredible, but as Marshall McLuhan said,"Only puny secrets need protection.
Big discoveries are protected by public incredulity."Now combine \textbf{the} electricity and \textbf{oil} revolutions, both driven by modern efficiency, and you get \textbf{the} really big story : reinventing fire, where business enabled and sped by smart policies in mindful markets can lead \textbf{the} United States completely off \textbf{oil} and \textbf{coal} by 2050, saving 5 trillion dollars, growing \textbf{the} economy 2. 6- fold, strengthening out national security, oh, and by \textbf{the} way, by getting rid of \textbf{the} \textbf{oil} and \textbf{coal}, reducing \textbf{the} \textbf{fossil} \textbf{carbon} \textbf{emissions} by 82 to 86 percent.
Now if you like any of those outcomes, you can support reinventing fire without needing to like all of them and without needing to agree about which of them is most important.
So focusing on outcomes, not motives, can turn gridlock and conflict into a unifying solution to America ' s energy challenge.
This also turns out to be \textbf{the} best way to cope with global challenges — climate change, nuclear proliferation, energy insecurity, energy poverty — all of which make us less safe.
Now our team at RMI helps smart companies to get unstuck and speed this journey via six sectoral initiatives, with some more hatching.
Of course there ' s still a lot of old thinking out there too.
Former \textbf{oil} man Maurice Strong said,"Not all \textbf{the} \textbf{fossils} are in \textbf{the} fuel."But as Edgar Woolard, who used to chair Dupont, reminds us,"Companies hampered by old thinking won ' t be a problem because,"he said,"they simply won ' t be around long-term."I ' ve described not just a once-in-a-civilization business opportunity, but one of \textbf{the} most profound transitions in \textbf{the} history of our species.
We humans are inventing a new fire, not dug from below, but flowing from above ; not scarce, but bountiful ; not local, but everywhere ; not transient, but permanent ; not costly, but free.
And but for a little transitional tail of natural gas and a bit of biofuel grown in ways that sustain and endure, this new fire is flameless.
Efficiently used, it really can do our work without working our undoing.
Each of you owns a piece of that \$5 trillion prize.
And our new book"Reinventing Fire"describes how you can capture it.
So with \textbf{the} conversation just begun at ReinventingFire. com, let me invite you each to engage with us and with each other, with everyone around you, to help make \textbf{the} world richer, fairer, cooler and safer by together reinventing fire.
Thank you.

% matched lemmas: airplane, backwards, buck, carbon, coal, computer, disappear, emission, fat, fossil, module, oil, one-third, pound, renewable, sum, technique, terrorist, textbook, the, transportation, two-third
\end{textsample}
